\section{Discussão}
\label{sec:discussao}

Propusemo-nos a medir se os LLMs respondem a perguntas de política de poluição
do ar com a mesma confiabilidade para o Sul Global e para o Norte Global, se o
enquadramento do prompt muda a resposta e que mecanismo impulsiona qualquer
lacuna. Em 25 países, 14 modelos, quatro idiomas, cinco tarefas e duas
personas, pontuados contra gabarito oficial por código onde a resposta é um
valor de registro e por um painel de três fornecedores nos demais casos, os
danos que conseguimos demonstrar são práticos: os três remédios a que um órgão
recorreria falham, o mais intuitivo deles torna a resposta materialmente pior,
e os modelos são menos confiáveis na recordação do valor vinculante em vigor.
Uma diferença de tier Norte/Sul persiste, e o mecanismo por trás dela não se
estabelece no nível em que a inferência é sólida.

\subsection{Três remédios intuitivos, todos falhando, e um que prejudica}

O idioma local não ajuda e prejudica materialmente. Perguntar no idioma do
próprio país reduz a acurácia em $\nnativa$~pp, significativamente nos três
idiomas e com o país como unidade de inferência, e a maior penalidade recai
sobre o hindi, o idioma de corpus mais escasso, com a Índia, o país de maior
pontuação entre os 25 em inglês, caindo mais. Essa é a direção que a
hierarquia de recursos prevê~\citep{joshi2020state,xue2021mt5}, e a medição
consegue vê-la porque a unidade é o idioma e não o país; com três idiomas, e
espanhol e português indistinguíveis, a ordenação é descritiva. O efeito não
depende de nenhum país, modelo, tarefa, condição de persona ou instrumento de
pontuação isolado, nem da qualidade da tradução, e se mantém em um desfecho que
não envolve juiz algum: o modelo fica marcadamente menos propenso a produzir um
número conferível, e mais propenso a não devolver nada.

A persona de gestor local não estreita a lacuna (DiD $\ndid$~pp, permutação
$p=\npersonap$, equivalente a zero em até $\pm2$~pp), em consonância com a
evidência de que personas de especialista não melhoram a acurácia
factual~\citep{zheng2024helpful,mollick2025personas}, e seu efeito principal é
uma pequena perda ($\npersonamain$~pp). Esse nulo pesa mais do que pesaria um
nulo sobre um hábito popular, porque o enquadramento de papel é a primeira
coisa que todo grande fornecedor diz aos usuários para fazer; Anthropic, OpenAI,
Google e xAI convergem no mesmo conselho (Seção~\ref{sec:related}), de modo que
uma administração que segue a documentação do fornecedor escreve a frase que
testamos. O que essa documentação reivindica para a frase é controle sobre tom,
formato e enquadramento, e nada nela diz que a técnica que molda o registro não
molda também a confiabilidade factual. Nossa estimativa separa as duas. Uma
checagem de manipulação confirma que metade das respostas assume
explicitamente o papel, de modo que o enquadramento é recebido e muda a voz;
ele não muda o que o modelo sabe. Testamo-lo digitado na mensagem e não
instalado como instrução de sistema, o canal que um servidor público numa
interface de chat tem (Seção~\ref{sec:limites}).

O modelo regional é o pior dos catorze ($\delta=\nregdelta$), e fica atrás do
seu par de escala nos prompts brasileiros em português para os quais foi
ajustado ($\nhthreebra$ contra $\nhthreebrallama$ do Llama~3.1~8B, em 20
células cada; $\nhthreenarrow$~pp nos 30 prompts em português), de modo que
recorrer a um modelo pequeno com ajuste local troca a escala e o treinamento de
fronteira que carregam a acurácia no domínio sem comprar de volta o
conhecimento local que promete. Três limites cercam isso: o teste estreito é
pequeno; só nos valores de registro, reservidos nos mesmos itens em português,
Cabra, seu modelo-base e o Llama~3.1~8B ficam a poucas respostas um do outro
(Seção~\ref{sec:res:robustez}), de modo que a margem do composto se apoia nas
tarefas e nos subcomponentes julgados; e o modelo que pudemos rodar é um ajuste
comunitário de uma base de 7B, não um modelo regional comercial de fronteira do
tipo que um órgão poderia adquirir. Para esses, o achado é uma previsão, e o protocolo liberado é o modo
de testá-la.

Os praticantes implantam os três como correções; nenhum funciona, e o que um
gestor mais provavelmente tenta primeiro causa dano ativo. Para órgãos
ambientais do Sul Global, nenhum atalho barato no nível do prompt ou do modelo
substitui a verificação do valor vinculante contra a fonte oficial.

\subsection{O piso de recordação está sobre os valores que vinculam}

O risco recai de forma desigual. Um gestor que redige prosa sobre política de
qualidade do ar é comparativamente bem servido; um gestor que recupera o ponto
de corte regulatório de que uma decisão depende está na pior célula do desenho
(Seção~\ref{sec:res:tarefa}), e pior ainda se perguntar num idioma de menor
recurso. O piso não é artefato de pontuação: o veredito factual nas duas
tarefas de recordação é decidido por código contra o registro, o contraste se
mantém com todo modelo e todo país como unidade, e um terço das respostas
erradas sobre o padrão nacional cita um valor da própria escada oficial do país
e não um valor inventado. Os modelos conhecem a escada e não sabem qual degrau
vincula.

\subsection{Uma diferença de tier que conseguimos mostrar, e um gradiente que não}

O achado geográfico é uma diferença entre grupos, não uma inclinação. Nos
desfechos que nenhum juiz toca, os países do Sul Global devolvem o valor do
registro $\ncodegap$~pp menos vezes do que os do Norte Global, com o país como
unidade; no composto a lacuna é de $\ntiergap$~pp, do mesmo tamanho dentro da
onda de extensão e nos quinze pré-especificados, e inalterada quando os quatro
membros da UE contam como um. Ela é concentrada: cresce com o quanto a resposta
depende de um fato específico do país e está ausente na única tarefa sem valor
de registro a acertar (Seção~\ref{sec:res:h1}), embora os intervalos por tarefa
com o país como unidade sejam largos e essa decomposição seja exploratória. O
gradiente de desenvolvimento fica abaixo do limiar pré-especificado em toda
leitura, e seu intervalo, $[\nrhocilo,\nrhocihi]$, inclui tanto o zero quanto o
limiar: 25 países nem estabelecem um gradiente nem o descartam. A ordenação dos
países mostra por que os dois resultados diferem. A Índia lidera os 25 países
apesar de ser do Sul Global, Canadá e França ficam baixos, de modo que a
desvantagem é uma diferença de tier sobreposta a idiossincrasia específica por
país que uma correlação de postos sobre 25 países não consegue atravessar.

\subsection{O mecanismo, e por que não se estabelece}

O teste de mecanismo pré-especificado falha, e parte da razão é um defeito no
instrumento. O proxy atribui o tamanho da edição da Wikipédia no idioma oficial
dominante do país, de modo que nove dos nossos 25 países recebem o valor da
Wikipédia inglesa. A variável não mede quanto texto existe no idioma de um país;
mede se o país tem inglês como língua oficial, o que acompanha a história
colonial e não o tamanho do corpus. Relatamos isso porque afeta qualquer estudo
que tenha usado ou venha a usar o mesmo proxy.

Mudar a unidade de variação deu uma associação no nível da resposta na direção
prevista, com as verificações que um efeito real passaria, e ela não sobrevive à
inferência com o país como unidade: agrupado por país, o coeficiente tem
$p=\nhfourclp$, os 25 déficits por país se correlacionam com a cobertura a
$\rho=\nhfourctryrho$ ($p=\nhfourctryp$), e dado o IDH a associação é
$\nhfourctrypartial$. Nesses 25 países, cobertura e desenvolvimento se
correlacionam a $\rho=\nrhocobhdi$, e o desenho não tem alavanca que mova um sem
o outro. Relatamos, portanto, a explicação por cobertura como compatível com os
dados e não estabelecida por eles, e a questão cobertura-ou-desenvolvimento
como aberta.

\subsection{A lição metodológica: a chave importa mais que a amostra}

As mesmas células pontuadas sustentam conclusões substantivas diferentes
conforme o instrumento usado para pontuá-las, e o tamanho dessa diferença é a
lição. Pontuar duas tarefas contra chaves provisórias e delegar a aritmética de
faixa a um juiz produziu um gradiente de desenvolvimento significativo e uma
penalidade de idioma pequena; reconstruir as chaves a partir de registros
oficiais e mover a comparação para código reduziu o gradiente pela metade e
mais que dobrou a penalidade (Seções~\ref{sec:res:h1} e \ref{sec:res:h2}).
Mover o veredito factual da terceira tarefa com valor de registro de dois
juízes de coleta para código mudou sua razão de chances-manchete em um terço,
porque os juízes estavam divididos pela mesma linha dos tiers. A própria chave
continuou se movendo sob escrutínio: a entrada do Reino Unido carregava uma
meta de 2040 no lugar do limite em vigor, e a de Bangladesh carregava o valor
de 2005 depois de as Rules de 2022 o terem elevado, de modo que um modelo que
citava o padrão vigente correto de um país do Sul Global era pontuado como
errado até que um leitor externo o apanhou. O conjunto de dados nunca mudou;
apenas o instrumento. O instinto da área quando um efeito de viés geográfico se
revela frágil é acrescentar países, mas o gabarito merece o escrutínio
primeiro. Uma chave defeituosa não apenas acrescenta ruído; acrescenta ruído
estruturado, porque os defeitos se concentram onde os textos oficiais são mais
difíceis de obter ou foram revisados mais recentemente, que é onde a hipótese
mora. O registro liberado carrega, por isso, um hash de conteúdo (material
suplementar), para que qualquer leitor saiba contra qual versão da chave um
número foi computado. A mesma lição vale para a inferência: com a resposta como
unidade, H4 tinha $p=\nhfourp$; com o país como unidade tem $p=\nhfourctryp$, e
só o segundo descreve o que 25 países conseguem mostrar.

\subsection{Por que o domínio importa}

O material particulado fino está entre os principais fatores de risco
ambiental para mortalidade prematura, e a carga recai de forma desproporcional
sobre o Sul Global: a Organização Mundial da Saúde atribui \num{88548} óbitos
no Brasil em 2021 à poluição do ar ambiente (intervalo de incerteza
\numrange{69477}{108431})~\citep{who2026gho_air41}, e a exposição de curto prazo
continua a produzir mortalidade mensurável nas cidades
brasileiras~\citep{requia2024shortterm}. Quando um gestor delega a um LLM uma
consulta normativa ou uma recomendação para um episódio, uma resposta imprecisa
pode se propagar para um parecer regulatório ou uma resposta de emergência.

Não afirmamos que o modelo é menos confiável onde as consequências são maiores,
porque não medimos isso: a carga de doença não está entre nossas covariáveis de
país, e a única observação no nível do país de que dispomos aponta na direção
oposta, já que a Índia carrega uma das maiores cargas de mortalidade por
PM\textsubscript{2,5} do mundo e é também o país de maior pontuação na nossa
amostra. Se a confiabilidade está desalinhada da necessidade é respondível com
a adição de uma covariável de mortalidade atribuível padronizada por idade, mas
não por este desenho. O arcabouço da colonialidade do
saber~\citep{mohamed2020decolonial} oferece uma interpretação do padrão que
medimos, em que os LLMs são resumos comprimidos de corpora produzidos sob
acesso digital e de publicação desigual e a penalidade de idioma é um elo entre
assimetria de corpus e assimetria de conhecimento; os dados sustentam a
penalidade e a diferença de tier, e a interpretação permanece uma
interpretação.

\subsection{Limitações e oportunidades futuras}
\label{sec:limites}

Cada limite abaixo cerca este desenho, e cada um nomeia o estudo que o
removeria. Nós os listamos na ordem em que os executaríamos.

\emph{Sem braço de recuperação.} Toda resposta vem da memória paramétrica do
modelo. Um gestor que usa um assistente com busca, ou um sistema de recuperação
sobre o diário oficial nacional, enfrenta outro instrumento, e a lacuna e a
penalidade de idioma podem mudar sob ele. Os prompts e registros liberados
permitem rodar um braço de recuperação sobre os mesmos itens, e é o primeiro
estudo que acrescentaríamos.

\emph{O enquadramento de papel foi testado no turno do usuário, não no prompt
de sistema.} O enquadramento de gestor é prefixado à mensagem em vez de
instalado como instrução de sistema, onde três dos quatro guias de fornecedor o
colocam; o da Google admite esse lugar ou o início do prompt do usuário, que é o
que usamos. A escolha segue a população que nos interessa, servidores públicos
na interface de chat, que não têm prompt de sistema a configurar, mas nosso
nulo cerca o canal desse usuário, não o de um desenvolvedor. Se uma persona em
instrução de sistema recupera acurácia factual onde uma digitada não recupera
é um experimento de duas células sobre os mesmos itens.

\emph{Do pontuado por painel ao ancorado em padrão-ouro.} Nenhum especialista
humano repontuou respostas aqui. Três propriedades cercam o que isso custa:
onde a resposta é um valor de registro, o veredito factual é computado por
código; onde é preciso julgamento, a pontuação é a média de três fornecedores
e não de um; e todo efeito relatado é um contraste entre grupos, em que a
severidade sistemática de um juiz se cancela. Restringir à acurácia factual
objetiva fortalece os dois efeitos-manchete, o oposto do que um juiz
permissivo produziria. Dois resíduos permanecem: os quatro subcomponentes não
factuais de uma célula de T1 ainda são os do juiz da coleta, razão pela qual o
veredito de código e não o composto é o desfecho primário de tier, e a
auditoria manual do extrator cobriu apenas T1. Uma validação cega por
especialistas humanos de uma amostra estratificada, cerca de 150 respostas
pontuadas por três avaliadores, que o pipeline liberado exporta diretamente,
converteria um resultado ancorado em fonte oficial num resultado ancorado em
padrão-ouro.

\emph{De um sinal dentro do país a um desenho causal.} Dentro dos países
removemos o confundimento no nível do país sem remover o confundimento
tarefa-por-país, e o sinal não alcança significância no nível do país. Um
desenho que varie a cobertura do país mantendo o desenvolvimento fixo, ou um
estudo longitudinal que acompanhe a acurácia à medida que a pegada
enciclopédica de um país cresce, trataria do primeiro problema; mais países
tratariam do segundo.

\emph{De um gradiente nulo a um teste adequadamente potente.} A diferença de
tier está estabelecida enquanto o gradiente não está, com $n=25$, e o intervalo
do gradiente inclui o limiar. Uma replicação pré-registrada com $n$ maior, que
o conjunto de dados aberto torna barata, decidiria se o gradiente está ausente
ou subpotente.

\emph{De três idiomas nativos à matriz multilíngue completa.} O resultado mais
forte se apoia em três idiomas nativos, de modo que seu mecanismo é relatado
descritivamente ainda que o efeito não seja; para um dos três, o hindi, o
próprio instrumento oficial é publicado em inglês, de modo que a penalidade
ali é uma penalidade pelo idioma do usuário e não pelo idioma do registro. A
objeção da qualidade da tradução está resolvida empiricamente pelo censo de
retrotradução relatado acima; o que permanece em aberto é a amplitude, e
escalar para mais idiomas é a extensão de maior valor que este estudo aponta.

\emph{De um domínio a um programa multidomínio.} Todos os efeitos vêm de um
domínio aplicado de alto risco e de uma família de corpus (Wikimedia), e
confinamos a linguagem mecanicista de acordo. A mesma suíte de tarefas e o
mesmo método de registro de gabarito se transferem diretamente a outros
domínios regulados e ancorados em fonte, como qualidade da água, segurança
alimentar e política de vacinação.

\emph{A pilha do modelo servido como unidade de análise.} Avaliamos modelos
tal como servidos por configurações de acesso fixas, que é o que um gestor
público consulta; uma lacuna residual poderia em princípio refletir a pilha em
vez dos pesos. A verificação de serviço do modelo regional
(Seção~\ref{sec:res:robustez}) cobre uma pilha; um desenho controlado com
pesos abertos idênticos em várias pilhas separaria as duas para todas elas, o
que a proveniência por chamada que liberamos permite.

\subsection{O que este estudo contribui}

As três correções a que os praticantes recorrem primeiro falham juntas: o
idioma local prejudica, o enquadramento de persona não faz nada, e o modelo de
7B com ajuste regional é o pior dos catorze. O resultado metodológico é que,
nesta literatura, a chave de gabarito, e o nível em que a inferência é feita,
importam mais que o tamanho da amostra. E o desenho separa o que 25 países
conseguem demonstrar do que não conseguem, relatando o segundo conjunto como
aberto e não como sustentado.
